% Source pp. 208-213 (Lectures19-25.pdf).
\section{Fixed Points and Criticality}

It is possible to establish direct relation between the singular behavior of
thermodynamic quantities near critical points and the fixed points of the
Wilson's RG. This relation explains experimentally observed \emph{universality}
of critical singularities.

Let us consider again a magnet. As we have seen few lectures ago, the long-range
fluctuations of local magnetization $M(x)$ can be taken into account by
considering the statistical integral
\[
Z=\int[D\varphi_0]\,e^{-\mathcal A_0[\varphi_0]}\,,
\]
where the field $\varphi_0(x)$ is proportional to $M(x)$,
\[
\varphi_0(x)=C(T)\,M(x)\,,
\]
with some temperature-dependent coefficient. It is assumed that this functional
integral has cutoff, with the cutoff parameter $\Lambda$ such that
\[
\Lambda^{-1}\gg a\,,
\]
the atomic scale. Exact form of the action $\mathcal A_0$ is not known; it
combines statistical contributions of all microscopic fluctuations and depends
significantly on all details of the microscopic interactions. In Landau-Ginzburg
approximation one takes
\[
\mathcal A_0=\int d^dx\,\left[\frac12\,(\partial\varphi_0)^2+\frac{m_0^2(T)}2\,\varphi_0^2+\frac{\lambda_0(T)}{4!}\,\varphi_0^4\right]\,,
\]
where $m_0^2(T)$, $\lambda_0(T)$ are unknown functions of $T$, but it can be any
other combination of powers of $\varphi_0$ and its derivatives. What is important,
all the coefficients are some (generally unknown) functions of the temperature
$T$. This means that the actions describing real physical system, with all its
microscopic details, being functions of the temperature, form one-dimensional
submanifold in $\Sigma$; I will call it the ``physical curve'' (Wilson used the
term ``canonical curve'').

Now let us assume that there is a fixed point $\mathcal A_*$ out there in
$\Sigma$. At this point we actually will consider a subspace
$\Sigma_{\text{symm}}\subset\Sigma$ which contains only actions invariant under
$\varphi_0\to-\varphi_0$. We assume that there is a fixed point
$\mathcal A_*\in\Sigma_{\text{symm}}$ such that it has only one (symmetric)
relevant (or marginally relevant) field $\Phi_0$, with dimension $D_0<d$. In
other words, the unstable manifold in $\Sigma_{\text{symm}}$ is one-dimensional,
and correspondingly the critical surface $\Sigma_c$ has codimension one. This
means that the physical curve generally intersects this critical surface at some
point. This is illustrated in figure~\ref{fig:physcurve}.
\begin{figure}[ht]
\centering\figfile{physcurve}
\caption{The critical surface $\Sigma_c$ with the fixed point $\mathcal A_*$ and
its one-dimensional unstable manifold $U(\mathcal A_*)$. The physical curve,
parametrized by the temperature, crosses $\Sigma_c$ at $T=T_c$.}
\label{fig:physcurve}
\end{figure}

Let us denote $T_c$ the temperature corresponding to the intersection point. The
large-distance properties of the system can be understood if we can follow the RG
flow which starts at any point of the physical curve. Let us consider a typical
pattern of the RG flow in this situation. Besides the critical fixed point
$\mathcal A_*$, there are two non-critical ($R_c=0$) fixed points,
$P_{\text{high}}$ and $P_{\text{low}}$\footnote{For a magnetic transition
$P_{\text{high}}$ corresponds to the probability distribution
$P_{\text{high}}[\varphi_0]\sim\prod_x\delta(\varphi_0(x))$, and
$P_{\text{low}}[\varphi_0]\sim\prod_x\delta(\varphi_0(x)-\bar\varphi_0)+\prod_x\delta(\varphi_0(x)+\bar\varphi_0)$.}
(figure~\ref{fig:critflow}).
\begin{figure}[ht]
\centering\figfile{critflow}
\caption{RG flow near a critical point. Trajectories starting on the physical
curve above $T_c$ go to $P_{\text{high}}$, those below $T_c$ go to
$P_{\text{low}}$. The trajectory $C$ starting at $T_c$ stays in $\Sigma_c$ and
ends at $\mathcal A_*$; a trajectory $C_\tau$ starting close to $T_c$ follows $C$,
lingers near $\mathcal A_*$, and leaves along the unstable manifold.}
\label{fig:critflow}
\end{figure}

A trajectory starting at any point on the physical curve corresponding to $T>T_c$,
and far enough from $T_c$, converges quickly to the noncritical fixed point
$P_{\text{high}}$. Similarly, starting at any point with $T<T_c$ one quickly
reaches $P_{\text{low}}$. The only trajectory which behaves qualitatively
differently is the one that starts at the point $T_c$ itself. This trajectory $C$
lies entirely at the critical surface, and converges to the critical fixed point
$\mathcal A_*$ as $l\to\infty$. This means that the large-distance physics exactly
at the critical temperature is completely determined by the fixed-point action
$\mathcal A_*$. In particular, the two-point correlation function will have the
scale-invariant form
\[
\bigl\langle M(x)M(x')\bigr\rangle_{T=T_c}=\frac\Const{|x-x'|^{2D_\varphi}}
\]
for $|x-x'|\gg a\,e^{l_0}$, where $l_0$ is some finite amount of the RG ``time''
$l$ needed to arrive at close proximity of $\mathcal A_*$. Here $D_\varphi$ is the
dimension of $\varphi$ at this fixed point.

Consider now a trajectory $C_\tau$ which starts at some point $T$ very close to
$T_c$,
\begin{equation}\label{eq:tau}
\tau=\frac{T_c-T}{T_c}\ll1\,.
\end{equation}
To be definite I will assume that $T<T_c$. There are three distinct stages of the
RG ``evolution'' for such trajectory. At the first stage it goes very close to $C$
approaching $\mathcal A_*$, until it comes to close proximity of this fixed point.
It then stays near $\mathcal A_*$ for very long RG ``time'' $l$. This is the
second stage. Finally, in the third stage it departs toward $P_{\text{low}}$
staying close to the unstable manifold $U$.

The second stage is the most important one; that is where the critical
singularities are formed. Because at this stage the trajectory stays close to the
fixed point, we can use linear analysis. Let us assume that in some finite amount
$l_0$ of RG ``time'' the trajectory $C_\tau$ comes close to $\mathcal A_*$, where
the corresponding action can be written as
\begin{equation}\label{eq:nearfp}
\mathcal A=\mathcal A_*+y_0\int d^dx\,\Phi_0(x)+\sum_{\alpha\neq0}y_\alpha\int d^dx\,\Phi_\alpha(x)\,,
\end{equation}
where the constants $y$ depend on $\tau$, the deviation \eqref{eq:tau} from the
critical point. It is important to note that this dependence is analytic because
it comes from integrating the RG flow equations for finite ``time'' $l_0$. As for
$\tau=0$ one must stay at the critical surface, which in the linear approximation
corresponds to $y_0=0$, for small $\tau$ we can assume
\begin{equation}\label{eq:y0T}
y_0=\Const\times(T_c-T)\,,
\end{equation}
and neglect the $\tau$-dependence of all $y_\alpha$ with $\alpha\neq0$.

In the second stage the RG flow is described by the equations
\[
\frac d{dl}\,y_0=\kappa_0\,y_0\,,\qquad\frac d{dl}\,y_\alpha=\kappa_\alpha\,y_\alpha\,,
\]
where the coefficients $\kappa$ are related to the dimensions $D_\alpha$ of the
fields $\Phi_\alpha$,
\[
\kappa_0=d-D_0>0\,,\qquad\kappa_\alpha=d-D_\alpha<0\quad\text{for}\quad\alpha\neq0\,.
\]
During this second stage all irrelevant terms in \eqref{eq:nearfp} die out, while
the coefficient $y_0$ increases as
\begin{equation}\label{eq:y0l}
y_0(l)=y_0\,e^{\kappa_0(l-l_0)}\,.
\end{equation}
Recall that if $R_c$ is the correlation length corresponding to the initial action
on the physical curve, the correlation length corresponding to the time $l$ along
the trajectory is
\[
R_c(l)=R_c\,e^{-l}\,.
\]
From this and \eqref{eq:y0l} one can find\footnote{Observe that within the linear
approximation \eqref{eq:nearfp} the ratio $R_c(l)/y_0^{-\nu}(l)$ is a constant
independent of $l$.}
\begin{equation}\label{eq:Rcnu}
R_c\sim y_0^{-\nu}\,,\qquad\text{where}\qquad\nu=\frac1{\kappa_0}=\frac1{d-D_0}\,,
\end{equation}
and from the linear dependence \eqref{eq:y0T}
\[
R_c\sim|T_c-T|^{-\nu}\,.
\]
Thus, the correlation length diverges near $T_c$, with the exponent determined by
the dimension $D_0$ of the relevant field $\Phi_0$. This divergence is in full
agreement with the idea that it is the large-scale fluctuations which are
responsible for the critical behavior.

Similar way one can relate other critical exponents to the dimensions of various
fields at the fixed point $\mathcal A_*$. For example, the spontaneous
magnetization
\[
\bar M\sim\langle\varphi_0\rangle
\]
(nonzero in the low-$T$ phase) must satisfy the CS equation, which in this case
reads
\[
\left[D_\varphi-\kappa_0\,y_0\,\frac\partial{\partial y_0}\right]\bar M=0\,,
\]
i.e.
\begin{equation}\label{eq:Mbar}
\bar M\sim y_0^{\frac{D_\varphi}{\kappa_0}}\,.
\end{equation}
Again, taking into account \eqref{eq:y0T} we find (for $T<T_c$)
\[
\bar M\sim(T_c-T)^\beta\,,\qquad\text{where}\qquad\beta=\frac{D_\varphi}{d-D_0}\,.
\]
It is useful to note that the relations \eqref{eq:Rcnu}, \eqref{eq:Mbar} can be
obtained by simple dimensional counting, if one assumes that the coupling constant
$y_0$ has the dimension $\kappa_0$ and the field $\varphi_0$ has the dimension
$D_\varphi$. This follows from the CS equation with the linear approximation for
the RG functions $\beta$; in fact it is valid beyond the linear approximation, as
it follows from the existence of ``canonical coordinates''.

As another example of this dimensional counting let us find some other critical
exponents in terms of dimensions. If, at $T=T_c$, one adds an external field,
i.e.\ the term
\[
H\int d^dx\,\varphi_0(x)\,,
\]
the dimensional counting implies that $H$ has dimension $d-D_\varphi$, and so the
magnetization as the function of $H$ must have the following singular behavior
\[
\bar M\sim H^{\frac1\delta}\,,\qquad\text{where}\qquad\delta=\frac{d-D_\varphi}{D_\varphi}\,.
\]
Similar way one finds for the specific heat at $H=0$
\[
C\sim|T_c-T|^{-\alpha}
\]
with
\begin{equation}\label{eq:alpha}
\alpha=2-d\,\nu\,.
\end{equation}
Note that all these exponents are expressed through just two numbers, the
dimensions of the fields $\varphi$ and $\Phi_0$. Therefore no matter what the
values of these dimensions are, the critical exponents must obey certain
\emph{scaling relations}, the relation \eqref{eq:alpha} being one example. These
relations are confirmed in experiment.

Note that the above analysis relies on the properties of the fixed point
$\mathcal A_*$, as well as on general topological properties of RG. Very little
indeed depends on the details of the original actions on the ``physical curve'',
and hence on the details of the microscopic structure of the magnet. One expects
therefore that the most important characteristics of the critical singularities
must not be different for phase transitions in different magnetic materials (with
the same symmetry). One fixed point may represent the criticality in the whole
class of physical systems. It is said it describes a \emph{universality class} of
critical behavior. In fact, systems of very different physical nature can belong
to the same universality class. For example, the liquid-gas critical point is
likely to represent the same universality class as the phase transition in axial
ferromagnet.

This all shows that understanding the fixed points is of crucial importance both
in statistical mechanics and in quantum field theory. Unfortunately, analysis and
classification of all fixed points are not yet achieved. In fact, at present very
few nontrivial fixed points can be analyzed on a systematic level.
